\begin{textsample}{POS Dim 9 – global\_north\_2024\_11 – Score 1.00 – global\_north\_2024\_11/118203359.txt}  \label{ex:f9_pos_402}
The PM expressed hope that this visit would help working groups follow up on the completion of these projects and enhance trade and economic cooperation between the two brotherly countries, noting that discussions also dealt with expanding and strengthening cooperation in humanitarian fields, especially given the deteriorating situation in Gaza and Sudan, with an agreement to work jointly on humanitarian support projects.
The PM confirmed that his meeting with the Egyptian FM addressed preparations for the humanitarian conference on Gaza, which is set to be held in Cairo on December 2.
He expressed the State of Qatar ' s eagerness to participate in the conference and provide all forms of support through it.
During the press conference, Sheikh Mohammed said that the meetings with Egyptian officials were an important opportunity to discuss the deteriorating regional situation, particularly with the ongoing war in Gaza and Lebanon, adding that warnings was issued of such conflicts and that it could potentially expand to encompass the entire region, stressing with regret that this is what is being witnessed today. @ @ @ @ @ @ @ @ @ @ consultation with its Egyptian brothers, particularly regarding negotiations aimed at achieving a ceasefire in Gaza as soon as possible.
He affirmed the necessity of coordinating with other brotherly nations on regional developments and ensuring a shared, clear vision among Arab countries in this regard.
For his part, the Minister of Foreign Affairs, Immigration and Egyptian Expatriates Affairs of the Arab Republic of Egypt welcomed the visit of Sheikh Mohammed to Cairo, stressing that it is a dear visit to the Egyptians.
He referred to the meeting of the PM with President Abdel Fattah El-Sisi, adding that it was a very friendly meeting, reflecting the warmth of relations between Egypt and Qatar The Egyptian FM explained that of Sheikh Mohammed met with HE Dr Mostafa Madbouly, the Egyptian Prime Minister, in the presence of a number of ministers from both sides.
The overall bilateral relations between the two countries were discussed, and the focus was on investment and trade aspects and enhancing trade exchange between the two countries.
He stressed that there @ @ @ @ @ @ @ @ @ @ countries to work on achieving a \textbf{qualitative} shift in relations between them in all fields, whether political, economic, cultural or commercial, and to enhance the existing consultations on various regional and international issues.
He pointed out that Sheikh Mohammed listened to a detailed presentation from the Egyptian Prime Minister and the ministers on the investment climate and business environment conducive to Qatari investments in Egypt, and working to overcome any difficulties that investors may face.
They also discussed trade exchange and cooperation in humanitarian fields."I was honored to meet with the Prime Minister and Minister of Foreign Affairs, and we discussed a number of regional and international files, most notably the continued Israeli aggression on the Gaza Strip and the violations in the West Bank.
We talked about the sincere and tireless Egyptian-Qatari efforts that have continued for more than a year to reach a deal that guarantees a ceasefire and stops the bloodshed of the brotherly Palestinian people, along with the release of all hostages and a number of Palestinian prisoners.
We @ @ @ @ @ @ @ @ @ @ consultation and coordination, as well as with sister Arab countries, so that there is common ground for the first Arab issue, which is the Palestinian issue,"he said.
The Egyptian FM stressed that the two sides discussed the great importance of the need for full and unconditional access to humanitarian aid to the Gaza Strip, adding that he informed the Prime Minister of all preparations related to the"Supporting Humanitarian Operations in Gaza"conference scheduled to be held in Cairo on December 2nd, which will be held under the patronage of HE President Abdel Fattah El-Sisi and Antonio Guterres, Secretary-General of the United Nations, which aims to enhance the volume of humanitarian aid in the Gaza Strip, deal with the current humanitarian disaster and provide the requirements for early recovery in the Strip.
He added,"I agreed with His Excellency the Prime Minister and Minister of Foreign Affairs that there is no way to achieve security and stability in the region except through the establishment of the Palestinian state on the entire @ @ @ @ @ @ @ @ @ @ and international humanitarian law, and respecting all relevant United Nations resolutions."Disclaimer : The content of this article is syndicated or provided to this website from an external third party provider.
We are not responsible for, and do not control, such external websites, entities, applications or media publishers.
The body of the text is provided on an"as is"and"as available"basis and has not been edited in any way.
Neither we nor our affiliates guarantee the accuracy of or endorse the views or opinions expressed in this article.
Read our full disclaimer policy here.

% matched lemmas: qualitative
\end{textsample}
